\clearpage
\thispagestyle{plain}
\markboth{PRELIMINARES}{PRELIMINARES}
\oldpage[IV-1]{5}
\phantomsection
\label{chapter:0IV1-es}

\begin{center}
{\large CAPÍTULO 0 (continuación)\footnotemark\par}
\vspace{1.5em}
{\Large\bfseries PRELIMINARES\par}
\vspace{2em}
{\large\bfseries Sumario\par}
\end{center}

\footnotetext{Recordemos que el símbolo $0_N$ (resp. $\operatorname{Err}_N$)
remite a la parte del capítulo 0 (resp. a la lista de Erratas) adjunta al
capítulo $N$.}

\begin{center}
\begin{tabular}{@{}r l@{}}
\S~14. & Dimensión combinatoria de un espacio topológico.\\
\S~15. & Sucesiones M-regulares y sucesiones $\mathscr F$-regulares.\\
\S~16. & Dimensión y profundidad en los anillos locales noetherianos.\\
\S~17. & Anillos regulares.\\
\S~18. & Complementos sobre las extensiones de álgebras.\\
\S~19. & Álgebras formalmente lisas y anillos de Cohen.\\
\S~20. & Derivaciones y diferenciales.\\
\S~21. & Diferenciales en los anillos de característica $p$.\\
\S~22. & Criterios diferenciales de lisitud formal y de regularidad.\\
\S~23. & Anillos japoneses.
\end{tabular}
\end{center}

Casi la totalidad de los parágrafos precedentes está dedicada a la exposición
de nociones de Álgebra conmutativa que se utilizarán a lo largo del capítulo
IV. Aunque buena parte de estas nociones figura ya en varias obras ([1], [12],
[13], [17], [30]), se ha considerado que sería más cómodo para el lector
disponer de una exposición continua y aproximadamente independiente. Junto con
los \S\S~5, 6 y 7 del capítulo IV (donde se utiliza el lenguaje de los
esquemas), estos parágrafos constituyen, dentro de nuestro Tratado, un pequeño
Tratado especial, aproximadamente independiente de los capítulos I a III, que
pretende presentar de forma coherente las propiedades de los anillos que «~se
comportan bien~» respecto de operaciones como la completación o la clausura
integral, vinculando sistemáticamente estas propiedades con algunas
concepciones generales\footnote{La mayoría de las propiedades de que se trata
fueron descubiertas por Chevalley, Zariski, Nagata y Serre. El método utilizado
aquí fue desarrollado por primera vez en el otoño de 1961, en un curso impartido
en Harvard University por A.~Grothendieck.}.

\setcounter{section}{13}
\section{Dimensión combinatoria de un espacio topológico}
\label{section:0.14-es}

\setcounter{subsection}{0}
\subsection{Dimensión combinatoria de un espacio topológico}
\label{subsection:0.14.1-es}

\oldpage[0\textsubscript{IV-1}]{6}
\begin{env}[14.1.1]
\label{0.14.1.1-es}
Sea $I$ un conjunto ordenado; una \emph{cadena} de elementos de $I$ es, por
definición, una sucesión finita estrictamente creciente
$i_0<i_1<\cdots<i_n$ de elementos de $I$ ($n\geq 0$); por definición, la
\emph{longitud} de esta cadena es $n$. Si $X$ es un espacio topológico, el
conjunto de sus subconjuntos \emph{cerrados irreducibles} está ordenado por
inclusión, de donde se obtiene la noción de \emph{cadena} de subconjuntos
cerrados irreducibles de $X$.
\end{env}

\begin{definition}[14.1.2]
\label{0.14.1.2-es}
Sea $X$ un espacio topológico. Se llama \emph{dimensión combinatoria de $X$}
(o simplemente \emph{dimensión de $X$}, si no puede haber confusión), y se
denota por $\dim_c(X)$, o simplemente por $\dim(X)$, a la cota superior de las
longitudes de las cadenas de subconjuntos cerrados irreducibles de $X$. Para
todo $x\in X$, se llama \emph{dimensión combinatoria de $X$ en $x$} (o
simplemente \emph{dimensión de $X$ en $x$}), y se denota por $\dim_x(X)$, al
número $\inf_U(\dim(U))$, donde $U$ recorre el conjunto de los entornos
abiertos de $x$ en $X$.
\end{definition}

De esta definición resulta que
\[
  \dim(\varnothing)=-\infty
\]
(la cota superior en $\overline{\mathbb R}$ del conjunto vacío es $-\infty$).
Si $(X_\alpha)$ es la familia de las componentes irreducibles de $X$, se tiene
\[
  \dim(X)=\sup_\alpha\dim(X_\alpha)
  \tag{14.1.2.1}\label{0.14.1.2.1-es}
\]
porque toda cadena de subconjuntos cerrados irreducibles de $X$ está, por
definición, contenida en una componente irreducible de $X$ y, recíprocamente,
estas componentes son cerradas en $X$, por lo que todo subconjunto cerrado
irreducible de un $X_\alpha$ es un subconjunto cerrado irreducible de $X$.

\begin{definition}[14.1.3]
\label{0.14.1.3-es}
Se dice que un espacio topológico $X$ es \emph{equidimensional} si todas sus
componentes irreducibles tienen la misma dimensión (igual entonces a $\dim(X)$
por (14.1.2.1)).
\end{definition}

\begin{proposition}[14.1.4]
\label{0.14.1.4-es}
\begin{enumerate}
  \item[\rm(i)] Para todo subconjunto $Y$ de un espacio topológico $X$, se tiene
  $\dim(Y)\leq\dim(X)$.
  \item[\rm(ii)] Si un espacio topológico $X$ es unión finita de cerrados
  $X_i$, se tiene $\dim(X)=\sup_i\dim(X_i)$.
\end{enumerate}
\end{proposition}

\begin{proof}
Para todo subconjunto cerrado irreducible $Z$ de $Y$, la adherencia $\overline Z$
de $Z$ en $X$ es irreducible ($0_{\mathrm I}$, 2.1.2) y
$\overline Z\cap Y=Z$, de donde se sigue (i). Por otro lado, si
$X=\bigcup_{i=1}^n X_i$, donde los $X_i$ son cerrados, todo subconjunto cerrado
irreducible de $X$ está contenido en uno de los $X_i$ ($0_{\mathrm I}$, 2.1.1),
y, por consiguiente, toda cadena de subconjuntos cerrados irreducibles de $X$
está contenida en uno de los $X_i$~; de donde se sigue (ii).
\end{proof}

De (14.1.4, (i)) se deduce que, para todo $x\in X$, también se puede escribir
\[
  \dim_x(X)=\lim_U\dim(U)
  \tag{14.1.4.1}\label{0.14.1.4.1-es}
\]
donde el límite se toma con respecto al conjunto filtrante decreciente de los
entornos abiertos de $x$ en $X$.

\oldpage[0\textsubscript{IV-1}]{7}
\begin{corollary}[14.1.5]
\label{0.14.1.5-es}
Sean $X$ un espacio topológico, $x$ un punto de $X$, $U$ un entorno de
$x$, e $Y_i$ ($1\leq i\leq n$) subconjuntos cerrados de $U$ tales que
$x\in Y_i$ para todo $i$ y que $U$ sea la reunión de los $Y_i$. Entonces se tiene
\[
  \dim_x(X)=\sup_i\bigl(\dim_x(Y_i)\bigr).
  \tag{14.1.5.1}\label{0.14.1.5.1-es}
\]
\end{corollary}

\begin{proof}
De (14.1.4, (ii)) resulta que
$\dim_x(X)=\inf_V(\sup_i(\dim(Y_i\cap V)))$ cuando $V$ recorre el conjunto
de los entornos abiertos de $x$ contenidos en $U$~; del mismo modo, se tiene
$\dim_x(Y_i)=\inf_V(\dim(Y_i\cap V))$ para todo $i$. El corolario es, pues,
evidente si
\[
  \sup_i\bigl(\dim_x(Y_i)\bigr)=+\infty~;
\]
en caso contrario, existe un entorno abierto $V_0\subset U$ de $x$ tal
que $\dim(Y_i\cap V)=\dim_x(Y_i)$ para $1\leq i\leq n$ y para todo
$V\subset V_0$, de donde se sigue la conclusión.
\end{proof}

\begin{proposition}[14.1.6]
\label{0.14.1.6-es}
Para todo espacio topológico $X$, se tiene
$\dim(X)=\sup_{x\in X}\dim_x(X)$.
\end{proposition}

\begin{proof}
De la definición (14.1.2) y de la prop. (14.1.4) resulta que
$\dim_x(X)\leq\dim(X)$ para todo $x\in X$. Por otro lado, sea
$Z_0\subset Z_1\subset\cdots\subset Z_n$ una cadena de subconjuntos cerrados
irreducibles de $X$, y sea $x\in Z_0$~; para todo abierto $U\subset X$
que contenga a $x$, $U\cap Z_i$ es irreducible ($0_{\mathrm I}$, 2.1.6) y
cerrado en $U$, y como se tiene $\overline{U\cap Z_i}=Z_i$ en $X$, los
$U\cap Z_i$ son distintos dos a dos~; por tanto, $\dim(U)\geq n$, lo que
termina la demostración.
\end{proof}

\begin{corollary}[14.1.7]
\label{0.14.1.7-es}
Si $(X_\alpha)$ es un recubrimiento abierto de $X$, o un recubrimiento cerrado
localmente finito de $X$, se tiene
$\dim(X)=\sup_\alpha(\dim(X_\alpha))$.
\end{corollary}

\begin{proof}
Si $X_\alpha$ es un entorno de $x\in X$, se tiene
$\dim_x(X)\leq\dim(X_\alpha)$, de donde se sigue la primera afirmación. Por otro
lado, si los $X_\alpha$ son cerrados y si $U$ es un entorno de $x\in X$ que
solo interseca un número finito de conjuntos $X_\alpha$, se tiene
\[
  \dim_x(X)\leq\dim(U)=\sup_\alpha\bigl(\dim(U\cap X_\alpha)\bigr)
  \leq\sup_\alpha\bigl(\dim(X_\alpha)\bigr)
\]
por (14.1.4), de donde se sigue la segunda afirmación.
\end{proof}

\begin{corollary}[14.1.8]
\label{0.14.1.8-es}
Sea $X$ un espacio de Kolmogórov ($0_{\mathrm I}$, 2.1.3) noetheriano, y
sea $F$ el conjunto de los puntos cerrados de $X$. Entonces
$\dim(X)=\sup_{x\in F}\dim_x(X)$.
\end{corollary}

\begin{proof}
Con la notación de la demostración de (14.1.6), basta observar
que existe en $Z_0$ un punto cerrado ($0_{\mathrm I}$, 2.1.3).
\end{proof}

\begin{proposition}[14.1.9]
\label{0.14.1.9-es}
Sea $X$ un espacio de Kolmogórov noetheriano no vacío. Para que
$\dim(X)=0$, es necesario y suficiente que $X$ sea finito y discreto.
\end{proposition}

\begin{proof}
Si un espacio $X$ es separado (y \emph{a fortiori} si $X$ es un espacio
discreto), los únicos subconjuntos cerrados irreducibles de $X$ son los
reducidos a un punto, por lo que $\dim(X)=0$. Recíprocamente, supongamos que
$X$ es un espacio de Kolmogórov noetheriano tal que $\dim(X)=0$~; como toda
componente irreducible de $X$ contiene un punto cerrado ($0_{\mathrm I}$, 2.1.3),
ha de reducirse a dicho punto. Como $X$ solo tiene un número finito de
componentes irreducibles, es finito y discreto.
\end{proof}

\begin{corollary}[14.1.10]
\label{0.14.1.10-es}
Sea $X$ un espacio de Kolmogórov noetheriano. Para que un punto $x\in X$ sea
aislado, es necesario y suficiente que $\dim_x(X)=0$.
\end{corollary}

\begin{proof}
La condición es evidentemente necesaria (sin hipótesis sobre $X$). También es
suficiente,
\oldpage[0\textsubscript{IV-1}]{8}
pues de ella resulta que $\dim(U)=0$ para un entorno abierto $U$ de $x$, y
como $U$ es un espacio de Kolmogórov noetheriano, $U$ es finito y discreto.
\end{proof}

\begin{proposition}[14.1.11]
\label{0.14.1.11-es}
La función $x\mapsto\dim_x(X)$ es semicontinua superiormente en $X$.
\end{proposition}

\begin{proof}
Es claro que esta función es semicontinua superiormente en todo
punto en el que vale $+\infty$. Supongamos, pues, que $\dim_x(X)=n<+\infty$~;
entonces, la fórmula (14.1.4.1) muestra que existe un entorno abierto $U_0$ de
$x$ tal que $\dim(U)=n$ para todo entorno abierto $U\subset U_0$ de $x$.
Dicho esto, para todo $y\in U_0$ y todo entorno abierto $V\subset U_0$ de $y$,
se tiene $\dim(V)\leq\dim(U_0)=n$ (14.1.4)~; se deduce, por tanto, de
(14.1.4.1) que $\dim_y(X)\leq n$.
\end{proof}

\begin{remark}[14.1.12]
\label{0.14.1.12-es}
Si $X$ e $Y$ son dos espacios topológicos, y $f:X\to Y$ una aplicación
continua, puede ocurrir que $\dim(f(X))>\dim(X)$~; se obtiene un ejemplo tomando
como $X$ un espacio discreto de dos elementos $a$, $b$, y como $Y$ el conjunto
$\{a,b\}$ dotado de la topología cuyos conjuntos cerrados son
$\varnothing$, $\{a\}$ y $\{a,b\}$~; si $f:X\to Y$ es la aplicación
identidad, se tiene $\dim(Y)=1$ y $\dim(X)=0$. Obsérvese que $Y$ es el espectro
de un anillo de valoración discreta $A$; su único punto cerrado es $a$ y su
punto genérico es $b$~; si $K$ y $k$ son, respectivamente, el cuerpo de
fracciones y el cuerpo residual de $A$, $X$ es el espectro del anillo $k\times K$
y $f$ la aplicación continua correspondiente al homomorfismo
$(\varphi,\psi):A\to k\times K$, donde $\varphi:A\to k$ y $\psi:A\to K$ son
los homomorfismos canónicos (cf. IV, 5.4.3).
\end{remark}

\subsection{Codimensión de una parte cerrada}
\label{subsection:0.14.2-es}

\begin{definition}[14.2.1]
\label{0.14.2.1-es}
Dada una parte cerrada irreducible $Y$ de un espacio topológico $X$, se llama
\emph{codimensión combinatoria} (o simplemente \emph{codimensión}) de $Y$ en
$X$, y se denota por $\operatorname{codim}(Y,X)$, al supremo de las longitudes
de las cadenas de partes cerradas irreducibles de $X$ cuyo menor elemento es
$Y$. Si $Y$ es una parte cerrada cualquiera de $X$, se llama
\emph{codimensión de $Y$ en $X$}, y se denota asimismo por
$\operatorname{codim}(Y,X)$, al ínfimo de las codimensiones en $X$ de las
componentes irreducibles de $Y$. Se dice que $X$ es
\emph{equicodimensional} si todas las partes cerradas irreducibles mínimas de
$X$ tienen la misma codimensión en $X$.
\end{definition}

De esta definición resulta que
$\operatorname{codim}(\varnothing,X)=+\infty$, pues el ínfimo de la parte vacía
de $\overline{\mathbb R}$ es $+\infty$. Si $Y$ es cerrada en $X$ y si
$(X_\alpha)$ (resp. $(Y_\beta)$) es la familia de las componentes irreducibles
de $X$ (resp. de $Y$), todo $Y_\beta$ está contenido en algún $X_\alpha$ y,
más generalmente, toda cadena de partes cerradas irreducibles de $X$ cuyo menor
elemento es $Y_\beta$ está formada por partes de algún $X_\alpha$; por tanto,
se tiene
\[
  \operatorname{codim}(Y,X)=
  \inf_\beta\left(\sup_\alpha
  \bigl(\operatorname{codim}(Y_\beta,X_\alpha)\bigr)\right)
  \tag{14.2.1.1}\label{0.14.2.1.1-es}
\]
donde, para cada $\beta$, $\alpha$ recorre el conjunto de los índices tales que
$Y_\beta\subset X_\alpha$.

\begin{proposition}[14.2.2]
\label{0.14.2.2-es}
Sea $X$ un espacio topológico.
\begin{enumerate}
  \item[\rm(i)] Si $\Phi$ es el conjunto de las partes cerradas irreducibles de
  $X$, se tiene
  \[
    \dim(X)=\sup_{Y\in\Phi}\bigl(\operatorname{codim}(Y,X)\bigr).
    \tag{14.2.2.1}\label{0.14.2.2.1-es}
  \]
  \oldpage[0\textsubscript{IV-1}]{9}
  \item[\rm(ii)] Para toda parte cerrada no vacía $Y$ de $X$, se tiene
  \[
    \dim(Y)+\operatorname{codim}(Y,X)\leq\dim(X).
    \tag{14.2.2.2}\label{0.14.2.2.2-es}
  \]
  \item[\rm(iii)] Si $Y$, $Z$ y $T$ son tres partes cerradas de $X$ tales que
  $Y\subset Z\subset T$, se tiene
  \[
    \operatorname{codim}(Y,Z)+\operatorname{codim}(Z,T)
    \leq\operatorname{codim}(Y,T).
    \tag{14.2.2.3}\label{0.14.2.2.3-es}
  \]
  \item[\rm(iv)] Para que una parte cerrada $Y$ de $X$ satisfaga
  $\operatorname{codim}(Y,X)=0$, es necesario y suficiente que $Y$ contenga
  una componente irreducible de $X$.
\end{enumerate}
\end{proposition}

\begin{proof}
Las afirmaciones (i) y (iv) son consecuencias inmediatas de la definición
(14.2.1). Para demostrar (ii), puede uno reducirse al caso en que $Y$ es
irreducible, y entonces la fórmula resulta de las definiciones (14.1.1) y
(14.2.1). Por último, para demostrar (iii), puede uno, en virtud de la
definición (14.2.1), reducirse primero al caso en que $Y$ es irreducible;
entonces $\operatorname{codim}(Y,Z)=
\sup_\alpha(\operatorname{codim}(Y,Z_\alpha))$ para las componentes
irreducibles $Z_\alpha$ de $Z$ que contienen $Y$; está claro que
$\operatorname{codim}(Y,T)\geq\operatorname{codim}(Y,Z)$, por lo que la
desigualdad es cierta si $\operatorname{codim}(Y,Z)=+\infty$; en caso
contrario, existe un $\alpha$ tal que
$\operatorname{codim}(Y,Z)=\operatorname{codim}(Y,Z_\alpha)$ y, en virtud de
(14.2.1), puede uno reducirse al caso en que $Z$ es también irreducible; pero
entonces la desigualdad (14.2.2.3) es una consecuencia evidente de la
definición (14.2.1).
\end{proof}

\begin{proposition}[14.2.3]
\label{0.14.2.3-es}
Sean $X$ un espacio topológico y $Y$ una parte cerrada de $X$. Para todo
abierto $U$ de $X$, se tiene
\[
  \operatorname{codim}(Y\cap U,U)\geq\operatorname{codim}(Y,X).
  \tag{14.2.3.1}\label{0.14.2.3.1-es}
\]
Además, para que los dos miembros de (14.2.3.1) sean iguales, es necesario y
suficiente que, si $(Y_\alpha)$ es la familia de las componentes irreducibles
de $Y$ que tienen intersección no vacía con $U$, se tenga
$\operatorname{codim}(Y,X)=
\inf_\alpha(\operatorname{codim}(Y_\alpha,X))$.
\end{proposition}

\begin{proof}
Se sabe ($0_{\mathrm I}$, 2.1.6) que $Z\mapsto\overline Z$ es una biyección
del conjunto de las partes cerradas irreducibles de $U$ sobre el conjunto de
las partes cerradas irreducibles de $X$ que tienen intersección no vacía con
$U$ y, en particular, hace corresponder a las componentes irreducibles de
$Y\cap U$ las componentes irreducibles de $Y$ que tienen intersección no vacía
con $U$; si $Y_\alpha$ es una de estas últimas,
se tiene, por tanto,
$\operatorname{codim}(Y_\alpha,X)=
\operatorname{codim}(Y_\alpha\cap U,U)$, y la proposición resulta entonces de
la definición (14.2.1).
\end{proof}

\begin{definition}[14.2.4]
\label{0.14.2.4-es}
Sean $X$ un espacio topológico, $Y$ una parte cerrada de $X$ y $x$ un punto de
$X$. Se llama \emph{codimensión de $Y$ en $X$ en el punto $x$}, y se denota
por $\operatorname{codim}_x(Y,X)$, al número
$\sup_U(\operatorname{codim}(Y\cap U,U))$, donde $U$ recorre el conjunto de
los entornos abiertos de $x$ en $X$.
\end{definition}

En virtud de (14.2.3), también puede escribirse
\[
  \operatorname{codim}_x(Y,X)=
  \lim_U\bigl(\operatorname{codim}(Y\cap U,U)\bigr)
  \tag{14.2.4.1}\label{0.14.2.4.1-es}
\]
donde el límite se toma sobre el conjunto filtrante decreciente de los
entornos abiertos de $x$ en $X$. Obsérvese que se tiene
\[
  \operatorname{codim}_x(Y,X)=+\infty\quad\text{si}\quad x\in X-Y.
\]

\oldpage[0\textsubscript{IV-1}]{10}
\begin{proposition}[14.2.5]
\label{0.14.2.5-es}
Si $(Y_i)_{1\leq i\leq n}$ es una familia finita de partes cerradas de un
espacio topológico $X$, y $Y$ es la unión de esta familia, se tiene
\[
  \operatorname{codim}(Y,X)=
  \inf_i\bigl(\operatorname{codim}(Y_i,X)\bigr).
  \tag{14.2.5.1}\label{0.14.2.5.1-es}
\]
\end{proposition}

\begin{proof}
En efecto, toda componente irreducible de uno de los $Y_i$ está contenida en
una componente irreducible de $Y$ y, recíprocamente, toda componente
irreducible de $Y$ es también componente irreducible de uno de los $Y_i$
($0_{\mathrm I}$, 2.1.1); la conclusión resulta, por tanto, de la definición
(14.2.1) y de la desigualdad (14.2.2.3).
\end{proof}

\begin{corollary}[14.2.6]
\label{0.14.2.6-es}
Sean $X$ un espacio topológico y $Y$ un subespacio cerrado localmente
noetheriano de $X$.
\begin{enumerate}
  \item[\rm(i)] Para todo $x\in X$, solo existe un número finito de
  componentes irreducibles $Y_i$ ($1\leq i\leq n$) de $Y$ que contienen $x$,
  y se tiene
  $\operatorname{codim}_x(Y,X)=
  \inf_i(\operatorname{codim}(Y_i,X))$.
  \item[\rm(ii)] La función
  $x\mapsto\operatorname{codim}_x(Y,X)$ es semicontinua inferiormente en $X$.
\end{enumerate}
\end{corollary}

\begin{proof}
En efecto, por hipótesis existe un entorno abierto $U_0$ de $x$ en $X$ tal
que $Y\cap U_0$ es noetheriano y, por tanto, solo tiene un número finito de
componentes irreducibles, que son las trazas sobre $U_0$ de componentes
irreducibles de $Y$; \emph{a fortiori}, solo existe un número finito de
componentes irreducibles $Y_i$ ($1\leq i\leq n$) de $Y$ que contienen $x$, y
se puede, sustituyendo $U_0$ por un entorno abierto $U\subset U_0$ de $x$ que
no tenga intersección con ninguno de los $Y_j$ que no contienen $x$, suponer
que los $Y_i\cap U$ son las componentes irreducibles de $Y\cap U$; para todo
abierto $V\subset U$ de $x$ en $X$, los $Y_i\cap V$ son entonces las
componentes irreducibles de $Y\cap V$, y (14.2.3) muestra entonces que
$\operatorname{codim}(Y_i,X)=\operatorname{codim}(Y_i\cap V,V)$, lo que
prueba (i). Además, para todo $x'\in U$, las componentes irreducibles de $Y$
que contienen $x'$ son algunas de las $Y_i$, por lo que
$\operatorname{codim}_{x'}(Y,X)\geq\operatorname{codim}_x(Y,X)$, lo que
demuestra (ii).
\end{proof}

\subsection{La condición de las cadenas}
\label{subsection:0.14.3-es}

\begin{env}[14.3.1]
\label{0.14.3.1-es}
En un espacio topológico $X$, diremos que una cadena
$Z_0\subset Z_1\subset\cdots\subset Z_n$ de subconjuntos cerrados irreducibles es
\emph{saturada} si no existe ningún subconjunto cerrado irreducible $Z'$ distinto
de los $Z_i$ y tal que $Z_k\subset Z'\subset Z_{k+1}$ para algún índice $k$.
\end{env}

\begin{proposition}[14.3.2]
\label{0.14.3.2-es}
Sea $X$ un espacio topológico tal que, para cualesquiera dos subconjuntos
cerrados irreducibles $Y$, $Z$ de $X$ tales que $Y\subset Z$, se tenga
$\operatorname{codim}(Y,Z)<+\infty$. Las dos condiciones siguientes son
equivalentes:
\begin{enumerate}
  \item[\rm(a)] Dos cadenas saturadas de subconjuntos cerrados irreducibles de $X$,
  con los mismos extremos, tienen la misma longitud.
  \item[\rm(b)] Si $Y$, $Z$, $T$ son tres subconjuntos cerrados irreducibles de
  $X$ tales que $Y\subset Z\subset T$, se tiene
  \[
    \operatorname{codim}(Y,T)=
    \operatorname{codim}(Y,Z)+\operatorname{codim}(Z,T).
    \tag{14.3.2.1}\label{0.14.3.2.1-es}
  \]
\end{enumerate}
\end{proposition}

\begin{proof}
Es inmediato que a) implica b). Recíprocamente, supongamos que se verifica b) y
demostremos que, si dos cadenas saturadas con los mismos extremos tienen
longitudes $m$ y $n\leq m$, se tiene necesariamente $m=n$. Razonemos por
inducción sobre $n$, siendo evidente la proposición para $n=1$. Supongamos,
pues, que $n>1$, $n<m$, y sea $Z_0\subset Z_1\subset\cdots\subset Z_n$ una
\oldpage[0\textsubscript{IV-1}]{11}
cadena saturada tal que existe otra cadena saturada con extremos
$Z_0$, $Z_n$ y de longitud $m$. Como
$\operatorname{codim}(Z_0,Z_n)\geq m>n$ y
$\operatorname{codim}(Z_0,Z_1)=1$, se deduce de b) que
$\operatorname{codim}(Z_1,Z_n)=\operatorname{codim}(Z_0,Z_n)-1>n-1$, lo que
contradice la hipótesis de inducción.
\end{proof}

Cuando se cumplen las condiciones de (14.3.2), se dice que $X$ satisface la
\emph{condición de las cadenas}, o también que es un \emph{espacio catenario}.
Es claro que todo subespacio cerrado de un espacio catenario es catenario.

\begin{proposition}[14.3.3]
\label{0.14.3.3-es}
Sea $X$ un espacio de Kolmogórov noetheriano de dimensión finita. Las
condiciones siguientes son equivalentes:
\begin{enumerate}
  \item[\rm(a)] Dos cadenas maximales de subconjuntos cerrados irreducibles de
  $X$ tienen la misma longitud.
  \item[\rm(b)] $X$ es equidimensional, equicodimensional y catenario.
  \item[\rm(c)] $X$ es equidimensional y, para dos subconjuntos cerrados
  irreducibles $Y$, $Z$ de $X$ tales que $Y\subset Z$, se tiene
  \[
    \dim(Z)=\dim(Y)+\operatorname{codim}(Y,Z).
    \tag{14.3.3.1}\label{0.14.3.3.1-es}
  \]
  \item[\rm(d)] $X$ es equicodimensional y, para dos subconjuntos cerrados
  irreducibles $Y$, $Z$ de $X$ tales que $Y\subset Z$, se tiene
  \[
    \operatorname{codim}(Y,X)=
    \operatorname{codim}(Y,Z)+\operatorname{codim}(Z,X).
    \tag{14.3.3.2}\label{0.14.3.3.2-es}
  \]
\end{enumerate}
\end{proposition}

\begin{proof}
Las hipótesis sobre $X$ implican que los extremos de una cadena maximal de
subconjuntos cerrados irreducibles de $X$ son necesariamente un punto cerrado y
una componente irreducible de $X$ ($0_{\mathrm I}$, 2.1.3); además, toda
cadena saturada con extremos $Y$, $Z$ (con $Y\subset Z$) está contenida en una
cadena maximal cuyos elementos distintos de los de la cadena dada o bien están
contenidos en $Y$, o bien contienen a $Z$. Estas observaciones establecen de
inmediato la equivalencia de a) y b), y prueban también que, si se verifica a),
se tiene, para todo subconjunto cerrado irreducible $Y$ de $X$,
\[
  \dim(Y)+\operatorname{codim}(Y,X)=\dim(X)~;
  \tag{14.3.3.3}\label{0.14.3.3.3-es}
\]
como se cumple (14.3.2.1), se deducen inmediatamente (14.3.3.1) y (14.3.3.2)
de (14.3.3.3). Recíprocamente, (14.3.3.1) implica (14.3.2.1), luego
(14.3.3.1) implica la condición de las cadenas en virtud de (14.3.2); además,
aplicando (14.3.3.1) al caso en que $Y$ se reduce a un punto cerrado $x$ de $X$
y $Z$ a una componente irreducible de $X$, se obtiene
$\operatorname{codim}(\{x\},X)=\dim(Z)$; se concluye que c) implica b).
Análogamente, (14.3.3.2) implica (14.3.2.1), luego la condición de las cadenas;
además, con la misma elección de $Y$ y $Z$ que antes, (14.3.3.2) implica
también $\operatorname{codim}(\{x\},X)=\dim(Z)$; por tanto (toda componente
irreducible de $X$ contiene un punto cerrado en virtud de ($0_{\mathrm I}$,
2.1.3)), d) implica b).
\end{proof}

Se dice que un espacio de Kolmogórov noetheriano es \emph{biequidimensional}
si es de \emph{dimensión finita} y verifica las condiciones equivalentes de
(14.3.3).

\begin{corollary}[14.3.4]
\label{0.14.3.4-es}
Sea $X$ un espacio de Kolmogórov noetheriano biequidimensional; entonces, para
todo punto cerrado $x$ de $X$ y toda componente irreducible $Z$ de $X$, se tiene
\[
  \dim(X)=\dim(Z)=\operatorname{codim}(\{x\},X)=\dim_x(X).
  \tag{14.3.4.1}\label{0.14.3.4.1-es}
\]
\end{corollary}

\oldpage[0\textsubscript{IV-1}]{12}
\begin{proof}
La última igualdad se deduce del hecho de que, si
$Y_0=\{x\}\subset Y_1\subset\cdots\subset Y_m$ es una cadena maximal de
subconjuntos cerrados irreducibles de $X$ y $U$ un entorno abierto de $x$, los
$U\cap Y_i$ son subconjuntos cerrados irreducibles de $U$ distintos dos a dos
(pues $\overline{U\cap Y_i}=Y_i$), de donde
$\dim(U)=\dim(X)$ en virtud de (14.1.4).
\end{proof}

\begin{corollary}[14.3.5]
\label{0.14.3.5-es}
Sea $X$ un espacio de Kolmogórov noetheriano; si $X$ es biequidimensional, lo
mismo ocurre con toda unión de componentes irreducibles de $X$ y con todo
subconjunto cerrado irreducible de $X$. Además, para todo subconjunto cerrado
$Y$ de $X$, se tiene entonces
\[
  \dim(Y)+\operatorname{codim}(Y,X)=\dim(X).
  \tag{14.3.5.1}\label{0.14.3.5.1-es}
\]
\end{corollary}

\begin{proof}
Como toda cadena de subconjuntos cerrados irreducibles de $X$ está contenida en
una componente irreducible de $X$, la primera afirmación se deduce
inmediatamente de (14.3.3). Además, si $X'$ es un subconjunto cerrado
irreducible de $X$, $X'$ verifica trivialmente las condiciones (14.3.3, c)), de
donde se sigue la segunda afirmación.

Por último, para demostrar (14.3.5.1), observemos que se ha visto en la
demostración de (14.3.3) que esta relación se verifica cuando $Y$ es
irreducible; si $Y_i$ ($1\leq i\leq m$) son las componentes irreducibles de
$Y$, aquella entre las $Y_i$ para la cual $\dim(Y_i)$ es mayor es también
aquella para la cual $\operatorname{codim}(Y_i,X)$ es menor; por tanto,
(14.3.5.1) se deduce de las definiciones de $\dim(Y)$ y de
$\operatorname{codim}(Y,X)$.
\end{proof}

\begin{remark}[14.3.6]
\label{0.14.3.6-es}
El lector observará que la demostración de (14.3.2) se aplica a un conjunto
ordenado cualquiera, pues el hecho de que se trate de subconjuntos cerrados
irreducibles de un espacio topológico no interviene en esta demostración. Lo
mismo ocurre con la demostración de (14.3.3) en un conjunto ordenado $E$ tal
que, para todo $x\in E$, existe un $z\leq x$ que es \emph{minimal} en $E$, y en
el que la longitud de las cadenas de elementos de $E$ está acotada.
\end{remark}
